\section{Related literature}\label{sec:literature}

Markowitz portfolio choice makes covariance the central input to
minimum-variance allocation \citep{markowitz_portfolio_1952}.
In short or high-dimensional return panels, however, the sample covariance
can be unstable or singular.
Regularization improves that input by replacing part of its sampling
variation with a structured target \citep{ledoit_wolf_2004}.
Such methods make covariance estimation more reliable, but they still solve
the portfolio problem by first estimating joint return risk.

The alternative developed here changes the source and form of the risk information.
Under a maintained information-to-exposure bridge, observable separation
between firms' information distributions rules out part of the
perfect-positive-dependence benchmark without estimating every covariance entry.
The resulting object is a one-sided bound that can be evaluated at feasible
portfolio weights and then minimized.
It is therefore a certificate about admissible joint risk configurations, not
a replacement point estimate for the covariance matrix.
This distinction determines how the factor, distributional, and textual
literatures enter the argument.

Factor models clarify the latent object that covariance summarizes.
The CAPM represents systematic covariance through scalar market loadings
\citep{sharpe_capital_1964}, whereas APT and approximate-factor models use
vector or expanding factor structures
\citep{ross_arbitrage_1976,chamberlain_rothschild_1983,bai_ng_2002}.
Characteristic-based models then use observable firm attributes to organize
loadings and expected returns
\citep{rosenberg1974extra,connor2007semiparametric,
connor2012efficient,kelly_characteristics_2019}.
These approaches explain how systematic exposures generate dependence, but an
observed characteristic is not itself an exposure or covariance estimate.

Distribution-valued characteristics replace a single firm descriptor with a
probability law and use quadratic transport to compare those laws.
For two exposure laws, the minimum expected squared displacement also
determines the maximum systematic covariance permitted by their marginals.
Using this identity, \citet{gawronsky_continuous_2026} derive a pairwise
covariance envelope from observable distributional separation under
maintained information-to-exposure restrictions.
That result establishes the pairwise level of the argument: observed geometry
can restrict how closely two latent systematic risks align.

At the cross-sectional level, \citet{gawronsky_spatial_2027} develop
target-anchored Wasserstein barycentric reconstruction.
The resulting distributional spanning weights form a target-specific barycentric
interaction field that enters exposure adjustment.
That cross-sectional system organizes firm-level exposure adjustment rather
than the risk of a portfolio assembled from those firms.
The two studies therefore supply neighbouring pairwise and cross-sectional
implications of distributional geometry, but neither resolves portfolio aggregation.

Portfolio risk adds a joint-compatibility requirement.
Pairwise optimal couplings need not be the pairwise marginals of any single
joint exposure law, so separately attainable covariance envelopes cannot
simply be stacked into a coherent portfolio risk configuration.
The portfolio-level step retains one joint law for all exposure marginals and
uses its weighted dispersion to deduct a certified amount from worst-case
systematic variance.
The weighted sum of squared Wasserstein distances provides a computational
relaxation of that multi-firm object, while the sharp form preserves the common
coupling needed for coherent portfolio risk.

Textual finance establishes that documents contain financially relevant information.
Prior studies map text into sentiment and disclosure measures, return
forecasts, predictive factors, and learned pricing objects
\citep{tetlock_giving_2007,loughran2011liability,
  gentzkow_text_as_data_2019,ke_predicting_returns_2019,
cong_textual_factors_2024,distaso_string_2024,wang2025newsnet}.
Their primary targets are expected returns, factors, or pricing rather than a
coherent bound on cross-asset portfolio risk.

Modern encoders make the distributional approach empirically feasible by
representing each document as a vector
\citep{reimers_sentencebert_2019,qwen3_embedding_2025}.
Word Mover's Distance provides an NLP precedent for using optimal transport
to compare empirical distributions of embeddings \citep{kusner_word_embeddings_2015}.
Retaining a firm's article embeddings as an empirical distribution preserves
within-firm heterogeneity that a single pooled vector would suppress.
In the present setting, those embeddings measure observed information
geometry; they do not measure latent exposure or covariance directly.
The carrier-and-slack restrictions provide the maintained link from that
measurement layer to exposure separation.

Taken together, portfolio theory supplies the decision problem, factor and
distributional models identify the latent risk objects and pairwise
restrictions, and text encoders supply observable information geometry.
What remains absent is a portfolio-level bridge from that geometry to a risk
bound supported by one coherent exposure law.
The theory therefore begins by separating observed characteristic laws,
latent exposure laws, and their maintained transmission before deriving the
coherent portfolio bound and its decision rule.
